\documentclass[prl,twocolumn]{revtex4-1}

\usepackage{graphicx}
\usepackage{color}
\usepackage{latexsym,amsmath}
\usepackage{amssymb}
\usepackage{subcaption}
\usepackage{comment}
\usepackage{float}
\definecolor{linkcolor}{rgb}{0,0,0.65} %hyperlink
\usepackage[pdftex,colorlinks=true, pdfstartview=FitV, linkcolor= linkcolor, citecolor= linkcolor, urlcolor= linkcolor, hyperindex=true,hyperfigures=true]{hyperref} %hyperlink%
\usepackage{enumitem}

\begin{document}

\title{Group 2607 -- Clustering}





\author{Fabio Cimino}
\author{Riccardo Ferrante}
\author{Daniele Ghezzi}
\author{Federico Scianna}

\date{\today}


\begin{abstract}
\noindent Clustering algorithms are widely used in unsupervised learning to identify hidden structures, detect anomalies, and group data according to intrinsic similarities.
However, their effectiveness strongly depends on the properties of the dataset under analysis.
In this work, we investigate the applicability and limitations of three widely used clustering techniques, K-means, DBSCAN, and hierarchical clustering, across four synthetic problems.
The considered scenarios include high-dimensional non-convex data, anomaly detection in autoregressive time series, and a checkerboard-shaped dataset designed to expose the limitations of standard clustering approaches.
Dimensionality reduction techniques, including PCA, t-SNE, and MDS, are employed to visualize complex data structures.
Clustering performance is evaluated through NMI, Silhouette Index, and F1 score.
The results show that clustering performance strongly depends on the geometric structure and density distribution of the data, highlighting the importance of selecting clustering techniques according to the intrinsic properties of the dataset.
\end{abstract}

\maketitle
\section{Introduction}
\noindent The aim of this assignment is to investigate the performance of different clustering techniques in a variety of unsupervised learning problems.
In particular, we analyze how the characteristics of a dataset influence the effectiveness of k-means and DBSCAN.
To this end, we consider four distinct problems that highlight the strengths and limitations of these clustering methods.\\
\noindent In the \textbf{first problem}, the objective is to identify clusters in a 
12-dimensional dataset.
Since direct visualization of such space is impossible, we first compare different dimensionality reduction techniques that allow for a 2-dimensional representation of the data.
We then evaluate the performance of several clustering algorithms. \\
\noindent The \textbf{second problem} focuses on anomaly detection in a collection of time series of equal length.
The goal is to isolate the anomalous series by identifying it as an outlier with respect to the remaining data, which are expected to form a single cluster. \\
\noindent In the \textbf{third problem}, we analyze another time series collection containing different hidden patterns embedded in noise.
The objective is to determine whether clustering techniques can distinguish time series containing sinusoidal or square-wave signals from those containing noise only.\\
\noindent Finally, in the \textbf{fourth problem}, we analyze a dataset that highlights the limits of the clustering techniques used.
\vspace{-1.5em}
\section{Data generation}
\noindent This section outlines the data generation process for the four problems.
Generated labels are retained solely for evaluating clustering performance.
\vspace{-1.5em}
\subsection{Problem 1}
\vspace{-1em}
\noindent For this experiment, we generate 600 3-dimensional data points, laying on different parametric curves: class 0 lies on a ring-like curve and is composed of 35\% of the data, class 1 on three translated ellipses and is composed of 45\% of the data, and class 2 on two symmetric oblique lines, and is composed of 20\% of the data. Gaussian noise is added to each point. 
The data is then embedded in a 12-dimensional space by appending nine additional noise dimensions.
Finally, the data is rotated with a random orthogonal matrix in order to hide the underlying structure. The original 3-dimensional data topology is shown in Figure~\ref{fig:placeholder}.
\begin{figure}[h]
    \centering
    \includegraphics[width=0.7\linewidth]{Pictures/experiment1/original data.png}
    \caption{3D data for Problem 1, before the addition of noise components to embed them in a 12D space.}
    \label{fig:placeholder}
\end{figure}\\
\vspace{-4em}
\subsection{Problem 2}
\vspace{-1em}
\noindent The data for this experiment are 101 time series of length 100. They are generated as autoregressive processes of first order AR(1):
\begin{equation}
    x_t = \phi \cdot x_{t-1} + \varepsilon_t, \ \varepsilon_t \sim \mathcal{N}(0, \sigma ^2),
\end{equation}
with a recall force parameter $\phi = 0.8$ and white noise standard deviation $\sigma = 5$. Every time series is interpreted as a point in $\mathbb{R}^{100}$. An anomaly is then introduced in one of the time series using the same process but adding a constant offset whose amplitude varies within $\{ 5,\ldots,25 \}$ with unitary increments.
\vspace{-1.5em}
\subsection{Problem 3}
\vspace{-1em}
\noindent The data for this problem consist of 120 time series generated as AR(1) processes. A third of the series are
augmented with a sinusoidal pattern, another third with
a square-wave pattern, while the remaining contain only
the AR(1) generated noise. The obtained set of time series is shown in Figure~\ref{fig:1}.
\vspace{-1.5em}
\subsection{Problem 4}
\vspace{-1em}
\noindent In this problem the data are 2-dimensional with binary labels, disposed in a checkerboard structure over a section of the plane as depicted in Figure~\ref{fig:2}.
\begin{figure}[b]
    \centering
    \includegraphics[width=\linewidth]{Pictures/experiment3/time_series.png}
    \caption{Dataset for problem 3: set of 120 time series with different anomalies.}
    \label{fig:1}
\end{figure} 
\begin{figure}[b]
    \centering
    \includegraphics[width=0.7\linewidth]{Pictures/experiment4/ground_truth.png}
    \caption{Labeled dataset for Problem 4.}
    \label{fig:2}
\end{figure} 

%\vspace{-2em}
\section{Methods}

\noindent In this section we illustrate the methods used for data visualization, clustering, and performance evaluation.

\vspace{-1.5em}
\subsection{Visualization Methods}
\vspace{-1em}
\noindent For data visualization, three dimensionality reduction methods are considered:
\vspace{-0.5em}
\begin{itemize}[leftmargin=0pt, itemsep=0pt]
    \item[] \textbf{Principal Component Analysis (PCA)}:
    a linear dimensionality reduction technique that projects data onto a new orthogonal basis formed by a subset of $k$ eigenvectors of the sample covariance matrix.
    The selected eigenvectors maximize the variance retained in the lower-dimensional representation.

    \item[] \textbf{t-distributed Stochastic Neighbor Embedding (t-SNE)}:
    a non-linear dimensionality reduction technique designed to preserve the local structure of high-dimensional data.
    The method first converts pairwise distances between data points in the original space into conditional similarity probabilities.
    For each point $x_i$, the similarity to another point $x_j$ is modeled as
    \begin{equation}
        p_{j|i} = \frac{\exp \left( -\|x_i - x_j\|^2 / 2\sigma_i^2 \right)}
        {\sum_{k \neq i} \exp \left( -\|x_i - x_k\|^2 / 2\sigma_i^2 \right)},
    \end{equation}
    where the bandwidth $\sigma_i$ is chosen such that the distribution has a prescribed perplexity, a parameter controlling the effective number of neighbors considered around each point.
    The conditional probabilities are then symmetrized to obtain joint similarities $p_{ij}$:
    \begin{equation}
        p_{ij} = \frac{p_{j|i} + p_{i|j}}{2N}.
    \end{equation}

    In the low-dimensional embedding, similarities between points $y_i$ and $y_j$ are instead modeled using a Student's t-distribution with one degree of freedom:
    \begin{equation}
        q_{ij} =
        \frac{\left(1+\|y_i-y_j\|^2\right)^{-1}}
        {\sum_{k \neq l}\left(1+\|y_k-y_l\|^2\right)^{-1}}.
    \end{equation}
    Compared to a Gaussian distribution, the heavier tails of the Student's t-distribution reduce the tendency of distant points to accumulate in the center of the embedding, allowing well-separated clusters to emerge more clearly.

    The embedding is obtained by minimizing the Kullback--Leibler divergence between the high-dimensional and low-dimensional similarity distributions:
    \begin{equation}
        C = \sum_{i \neq j} p_{ij}\log\frac{p_{ij}}{q_{ij}},
    \end{equation}
    which is optimized through gradient descent.
    As a result, points that are close in the original space tend to remain close in the low-dimensional representation, while dissimilar points are mapped farther apart.

    \item[] \textbf{Metric Multidimensional Scaling (Metric MDS)}:
    a geometry-based technique that seeks a low-dimensional embedding preserving global pairwise distances.
    Given a distance matrix $d_{ij}$ in the original space, Metric MDS learns low-dimensional coordinates $y_i$ by minimizing the stress function:
    \begin{equation}
        S = \sum_{i<j} \left( d_{ij} - \|y_i - y_j\| \right)^2,
    \end{equation}
    ensuring that the spatial relationships in the lower-dimensional representation approximate the original distances. The algorithm takes as input the number of components of the low-dimensional space. It also depends on two other parameters, $n_{init}$ and $max_{iter}$. The former is the number of times the algorithms restarts, while the latter is the maximum number of iterations after which the optimization process stops.

    
\end{itemize}


\vspace{-1.5em}
\subsection{Clustering Algorithms}
\vspace{-1em}
\noindent Clustering aims to partition a dataset into groups of similar points based on a defined notion of similarity.
In this work, three widely used clustering algorithms are considered:
\vspace{-0.5em}
\begin{itemize}[leftmargin=0pt, itemsep=0pt]
    \item[] \textbf{K-means}:
    a centroid-based algorithm that partitions the dataset into $K$ clusters by minimizing the cost function
    \begin{equation}
        G=\sum_{j=1}^{K}\sum_{x\in C_j}||x-\mu_j||^2.
    \end{equation}
    The algorithm iteratively updates both the clusters $C_j$ and their centroids $\mu_j$ for all $j\in\{1,\ldots,k\}$ until convergence.

    \item[] \textbf{Density-Based Spatial Clustering of Applications with Noise (DBSCAN)}:
    groups points based on local density.
    Points that are close to many neighbors are grouped together, while points in sparse regions are treated as outliers.
    This approach allows DBSCAN to identify clusters of arbitrary shape and to automatically detect outliers, without requiring the number of clusters as an input. The algorithm depends on two parameters: $\varepsilon$ and $m_s$. The former is the maximum distance between two points for them to be considered neighbors, while the latter represents the minimum number of neighbors within a radius $\varepsilon$ needed to define a cluster.
    
    
\begin{comment}
    \item \textbf{Spectral Clustering}
    is a graph-based technique capable of identifying non-convex clusters by analyzing the connectivity of the data.
    An affinity matrix is first constructed to model pairwise similarities, here using a Radial Basis Function (RBF) kernel where the hyperparameter $\gamma$ acts as the inverse of the variance, controlling the scale of local neighborhoods.
    A diagonal degree matrix is computed by summing the rows of the affinity matrix, and the Graph Laplacian is obtained by subtracting the affinity matrix from the degree matrix.
    The eigenvectors corresponding to the smallest eigenvalues of the Laplacian are used to project the data into a lower-dimensional spectral space, where connected components become separable.
    Finally, a QR decomposition is applied to extract discrete cluster labels.
\end{comment}

\item[] \textbf{Hierarchical Clustering}:
    builds a nested sequence of clusters by iteratively merging smaller groups based on a defined linkage criterion. The resulting hierarchical structure is visualized using a \textit{dendrogram}, a tree-like diagram that records the distance or dissimilarity at which each merge occurs. This approach allows for the determination of $K$ by drawing a horizontal cut across the dendrogram at a specific height, defining the clusters based on a desired distance threshold.

\end{itemize}

\vspace{-1.5em}
\subsection{Performance Metrics}
\vspace{-1em}
\noindent Finally, we implement several metrics to assess the quality of the clustering:
\vspace{-0.5em}
\begin{itemize}[leftmargin=0pt, itemsep=0pt]
    \item[] \textbf{Normalized Mutual Information (NMI)}:
    used to quantify the information shared between the clusters identified by the algorithm and the ground truth labels provided in the dataset.
    It can be evaluated as 
    \begin{equation}
        \mathrm{NMI}(Y,C)=\frac{2I(Y;C)}{H(Y)+H(C)},
    \end{equation}
    where $I$ represents the mutual information between real labels $Y$ and predicted ones $C$. The normalization factor is given by the sum of the entropies $H$ associated to both $Y$ and $C$ \cite{vinh2010information}.
    The notion of entropy adopted is the information entropy, where the probability is calculated as the ratio of labels of a certain class $i$ over all the total number of labels. \\
    NMI values range from zero to one, with the latter indicating that the clusters perfectly correspond to the known classes and the former indicating no correlation.
    
    \item[] \textbf{Silhouette Index}:
    evaluates how well each point is assigned to its cluster by comparing the average distance to other points in the same cluster with the average distance to points in the nearest neighboring cluster.
    For each data point, the silhouette value is calculated as 
    \begin{equation}
        s = \frac{b - a}{\max(a, b)},
    \end{equation}
    where $a$ is the average intra-cluster distance and $b$ is the smallest average inter-cluster distance to any other cluster.
    The overall Silhouette Index is the mean silhouette value across all points.
	Values close to one indicate that points are well matched to their own clusters and poorly matched to neighboring clusters, values near zero indicate overlapping clusters, while negative values indicate possible misclassification.
    
    \item[] \textbf{F1 Score}:
    provides a combined measure of the precision and recall of the clustering algorithm with respect to the ground truth reference class.
    Precision measures the fraction of points assigned to a cluster that actually belong to the corresponding class. It can be computed as
    \begin{equation}
        \text{Precision} = \frac{TP}{TP+FP},
    \end{equation}
    where $TP$ and $FP$ respectively represent \textit{True Positives} and \textit{False Positives}.
    Recall measures the fraction of points of the same class correctly captured within a cluster. Recall can be derived as
    \begin{equation}
        \text{Recall} = \frac{TP}{TP+FN},
    \end{equation}
    where $FN$ refers to \textit{False Negatives}.
    The F1 score is the harmonic mean of precision and recall.
    High values indicate that clusters accurately represent the true classes, balancing the trade-off between including only relevant points and capturing all members of a ground truth class.
    
\end{itemize}

\noindent Algorithms and metrics are implemented using the Python package \texttt{Scikit-Learn} \cite{geron}.

\begin{figure*}[!t]
    \begin{subfigure}{0.4\linewidth}
        \centering
        \includegraphics[width=\linewidth]{Pictures/experiment1/t-SNE.png}
        \caption{t-SNE visualization}
        \label{fig:visualization_tsne}
    \end{subfigure}
    \hspace{1.5cm}
    \begin{subfigure}{0.4\linewidth}
        \centering
        \includegraphics[width=\linewidth]{Pictures/experiment1/MDS.png}
        \caption{MDS visualization}
        \label{fig:visualization_mds}
    \end{subfigure}
    \caption{Comparison of t-SNE and MDS reduction methods. The plots are colored according to the true labels.}
    \label{fig:visualization}
\end{figure*}
\newpage
\section{Results}
\subsection{Problem 1}
\vspace{-1em}
\noindent Before running any algorithm, we used different visualization techniques in order to retrieve the original 3D structure of the data. First, we used PCA to reduce the data to the three principal components. The resulting plots, colored using the true labels, help to visualize the original topology. However, in order to have a clear understanding of the overall structure of the dataset, three plots are needed, corresponding to the projections of the data on the $PC_1-PC_2$, $PC_2-PC_3$, $PC_1-PC_3$ planes. For this reason, more compact visualization methods such as t-SNE and MDS are used. \\
\noindent For t-SNE, a search for the optimal perplexity parameter was conducted over the range $[5,100]$ with increments of 5, using the silhouette score as the evaluation metric.
The search identified the optimal perplexity value as $p_{best}=40$, which was then used to generate the two-dimensional visualization presented in Figure~\ref{fig:visualization_tsne}.
For MDS, the number of output dimensions was set to $n_d = 2$ to produce a two-dimensional plot.
A small number of initializations $n_{init}=2$ was chosen due to the linear scaling of run time with this parameter and the negligible improvement in visual quality observed for higher values.
The maximum number of iterations was set to $max_{iter} = 200$.\\
\noindent Following the evaluation of both t-SNE and MDS, as illustrated in Figure~\ref{fig:visualization}, MDS was selected for the presentation of the results.
This choice was motivated by the method’s ability, similarly to PCA, to preserve the global structure and overall patterns of the original dataset while providing a more compact representation.
In contrast, although t-SNE offers a more compact visualization than PCA, it introduces distortions in the global data structure, making it less suitable for the present analysis.
\noindent Having established a suitable visualization framework, clustering algorithms were subsequently applied.
The first algorithm considered was K-means.
To determine the optimal number of clusters K, hierarchical clustering was performed, and the optimal value was selected based on the largest jump in the dendrogram, as shown in Figure~\ref{fig:dendrogram}.
\begin{figure}[b]
    \centering
    \includegraphics[width=\linewidth]{Pictures/experiment1/dendrogram.png}
    \caption{The dendrogram resulting from hierarchical clustering in Problem 1.}
    \label{fig:dendrogram}
\end{figure}
\noindent The best value for the $K$ parameter was found to be $K=2$. The results obtained by initializing the algorithm with this value are reported in Figure~\ref{fig:kmeans}. 
\begin{figure}[t]
    \centering
    \includegraphics[width=0.9\linewidth]{Pictures/experiment1/K_means.png}
    \caption{Clusters obtained from the K-means algorithm for Problem 1.}
    \label{fig:kmeans}
\end{figure}
It is evident that K-means failed to capture the underlying structure of the dataset.
This outcome was expected, as K-means assumes convex, linearly separable clusters, an assumption violated by the present dataset.
The inadequacy of this clustering is further confirmed by the NMI score of $0.0105$, which is close to zero.
\noindent Next, the DBSCAN algorithm was applied. The minimum sample parameter $m_s$ was set to 3, as $m_s=1$ would reduce each data point to a core point, eliminating the density requirement, and $m_s=2$ effectively reduces DBSCAN to hierarchical clustering with single linkage.
The $\varepsilon$ parameter was estimated by identifying the point of maximum second derivative of the $k$-NN distance curve, after applying Savitzky--Golay smoothing~\cite{savitzky1964smoothing}, as illustrated in Figure~\ref{fig:knn}.

\noindent The resulting parameters were
\begin{equation}
    (\varepsilon, m_s) = (249.01, 3).
\end{equation}
\noindent The obtained value of $\varepsilon$ is consistent with the result of a grid search over the NMI, which yielded $\varepsilon_{\text{grid-search}}=255.0$, as shown in Figure~\ref{fig:nmi_heatmap}.
\begin{figure}[b]
    \centering
    \includegraphics[width=0.85\linewidth]{Pictures/experiment1/Heatmap_NMI.png}
    \caption{Grid-search for $(\varepsilon, m_s)$ parameters in DBSCAN for Problem 1.}
    \label{fig:nmi_heatmap}
\end{figure}
\noindent Figure~\ref{fig:dbscan} presents the clustering results obtained with DBSCAN.
While the algorithm successfully identifies the cluster corresponding to the ring-like curve, it merges the remaining two clusters, likely prioritizing density over geometric shape.
Additionally, four points are identified as outliers.
The resulting NMI score of $0.7527$ represents a substantial improvement over K-means, whose clusters failed to capture the dataset’s structure.
Nevertheless, this is still not a satisfying result, since the algorithm failed to correctly recognise all three groups lying in the data.
\begin{figure}[t]
    \centering
    \includegraphics[width=\linewidth]{Pictures/experiment1/k_nn.png}
    \caption{$k$-NN distance plot for $m_s=3$ in Problem 1.}
    \label{fig:knn}
\end{figure}  
\begin{figure}[b]
    \centering
    \includegraphics[width=0.9\linewidth]{Pictures/experiment1/DBSCAN.png}
    \caption{Clusters obtained using the DBSCAN algorithm for problem 1.}
    \label{fig:dbscan}
\end{figure}

\subsection{Problem 2}
\begin{figure*}[!t]
    \centering
    \includegraphics[width=0.8\linewidth]{Pictures/experiment2/Anomalies_amplitude13.png}
    \hskip 1mm
    \includegraphics[width=0.8\linewidth]{Pictures/experiment2/Anomalies_amplitude15.png}
    \caption{Ground truth anomalies for $A=13$ (top panel) and $A=15$ (bottom panel) and the correspondent results of DBSCAN.}
    \label{fig:Anomalies_amplitude}
\end{figure*}

\noindent To address Problem 2, DBSCAN was preferred over K-means due to the intrinsic limitations of centroid-based clustering in anomaly detection. Specifically, K-means partitions the dataset by minimizing the within-cluster variance. In this case, K-means has to be initialized with $K=2$ in order to collect every non anomalous time series in one cluster, while keeping the anomaly in the other. If one centroid were to correctly isolate the anomaly, the second centroid would have to cover the entire diffuse cloud of normal time series, resulting in a disproportionately high cost function. To minimize the overall variance, K-means tends to arbitrarily bisect the dataset into two roughly equal-variance clusters, inevitably grouping the anomaly together with normal time series.
\begin{figure}[b]
    \centering
    \includegraphics[width=0.7\linewidth]{Pictures/experiment2/DBSCAN_performance.png}
    \caption{F1 score of DBSCAN for different offset amplitudes of the injected anomaly.}
    \label{fig:DBSCAN_performance}
\end{figure}

\noindent In contrast, DBSCAN effectively isolates the anomaly by evaluating Euclidean distances between time series in the 100-dimensional space.
The algorithm requires the specification of $\varepsilon$ and $m_s$ parameters.
For $m_s=1$, the algorithm would have classified the anomaly as part of the cluster containing every other non-anomalous time series, preventing effective anomaly detection.
Any value greater than 1 is instead appropriate. \\
To estimate the optimal $\varepsilon$, the pairwise distances between all time series were analyzed. The squared Euclidean distance between two time series $X$ and $Y$ is given by
\begin{equation}
    D^2=\sum_{t=1}^L(X_t-Y_t)^2.
\end{equation} 
The average distance between time series was then obtained by evaluating the mean of this distribution, yielding a value of $d_{mean}\approx118$.
This empirical finding is validated by the theoretical expectation for independent AR(1) processes \cite{hamilton1994time}. Since the stationary variance of such processes is 
\begin{equation}
\sigma_t^2 = \frac{\sigma^2}{1-\phi^2}.
\end{equation}
The expected Euclidean distance between two regular trajectories is given by
\begin{equation}
    d_{mean}=\sigma\sqrt{\frac{2L}{1-\phi^2}}.
\end{equation}
We retrieve a theoretical distance in agreement with the empirical estimation.
The final value of $\varepsilon$ was chosen by adding a 10\% margin to $d_{mean}$ to ensure that the anomaly would be classified as an outlier.
Smaller values of $\varepsilon$ risk detecting false anomalies, whereas larger values would fail to identify the anomaly correctly.\\
\noindent The algorithm was then executed on the dataset while varying the amplitude of the anomalous offset.
Figure~\ref{fig:DBSCAN_performance} presents the F1 score as a function of the anomaly amplitude and two operational regimes are evident from the results.
For amplitudes $A<15$, DBSCAN is unable to distinguish between the anomalous and regular time series.
For $A\ge 15$, the algorithm successfully identifies the anomalous series. Figure~\ref{fig:Anomalies_amplitude} illustrates the ground truth anomalies and the corresponding DBSCAN results for two representative amplitudes.\\
\noindent In each subfigure, the left panel displays the injected anomaly in isolation for visual comparison, while the right panel shows the full dataset labeled according to DBSCAN clustering.
These findings generalize to multiple anomalies; however, a single anomaly was injected here to maintain graphical clarity, facilitating interpretation of the results.
\begin{figure}[t]
    \centering
    \includegraphics[width=\linewidth]{Pictures/experiment3/dendrogram.png}
    \caption{The dendrogram resulting from hierarchical clustering in Problem 3.}
    \label{fig:3}
\end{figure}
\begin{figure}[b]
    \centering
    \includegraphics[width=\linewidth]{Pictures/experiment3/kmeans_grid_search.png}
    \caption{Grid-search of the parameter K in Problem 3.}
    \label{fig:4}
\end{figure}
\vspace{-2em}
\subsection{Problem 3}
\noindent In order to correctly use K-means algorithm to solve this problem, the best $K$ parameter is determined. To do so, hierarchical clustering is performed, selecting the optimal value for $K$ from the maximum jump that occurs in the dendrogram, reported in Figure~\ref{fig:3}. This analysis yielded a value of $K = 3$. This value was confirmed by a parameter search using the NMI as an evaluation metric. The obtained maximum value is
\begin{equation}
    \max_K \mathrm{NMI} = \mathrm{NMI}(K = 3) \approx 0.90,
\end{equation}
as shown in Figure~\ref{fig:4}. It is important to note that this second method makes use of the ground truth labels. Therefore it does not constitute an estimation of $K$, but rather a validation of the result obtained from the dendrogram.
\newline K-means algorithm was then executed setting $K=3$. The results are reported in Figures~\ref{fig:5}-\ref{fig:6}.
\begin{figure}[t]
    \centering
    \includegraphics[width=\linewidth]{Pictures/experiment3/clusters.png}
    \caption{Clusters obtained from K-means algorithm in shape-based clustering in Problem 3.}
    \label{fig:5}
\end{figure}
\begin{figure}[b]
    \centering
    \includegraphics[width=\linewidth]{Pictures/experiment3/centroids.png}
    \caption{Centroids obtained from K-means algorithm in shape-based clustering in Problem 3.}
    \label{fig:6}
\end{figure}
The centroids reported in Figure~\ref{fig:6} confirm that K-means has successfully recovered the underlying patterns: centroid 0 resembles a square wave, centroid 1 a  sinusoidal wave, centroid 2 a flat, noisy signal. Thus the three groups defined in the data generation process have been recovered. The resulting clusters are depicted in Figure~\ref{fig:6}.
\begin{figure}[t]
    \centering
    \includegraphics[width=0.6\linewidth]{Pictures/experiment4/k_means.png}
    \caption{Clusters obtained from K-means algorithm in Problem 4. DA CAMBIARE}
    \label{fig:7}
\end{figure}

\subsection{Problem 4}
\noindent To address Problem 4, K-means and DBSCAN were applied to recover the original labels. However, both algorithms proved unsuccessful. K-means failed due to the non-linear separability of the data, since it assumes convex, linearly separable clusters, which does not hold true for this dataset. Conversely, DBSCAN struggled with the dataset's uniform density. As a density-based approach, it merged all data points into a single cluster, withoud detecting any internal density variations.
The results are reported in Figure~\ref{}.


\begin{figure}[H]
    \centering
    \includegraphics[width=0.6\linewidth]{Pictures/experiment4/DBSCAN.png}
    \caption{Clusters obtained from DBSCAN algorithm in Problem 4. DA CAMBIARE}
    \label{fig:8}
\end{figure}


\section{Conclusions}
\noindent In this work, we evaluated the performance of K-means, DBSCAN, and hierarchical clustering across four synthetic datasets, each designed to test different aspects of clustering algorithms.
K-means ....
DBSCAn ...
Hierarchical clustering ...
No clustering algorithm proved effective in recovering the underlying structure of a dataset with uniform density and non-linearly separable pattern.
In conclusion the results demonstrate that the effectiveness of clustering algorithms depends on the dataset’s intrinsic properties, and that combining visual representations with quantitative metrics is essential for robust evaluation.


%-------BIBLIOGRAPHY-------------------

\nocite{*} 
\bibliographystyle{apsrev4-1}
\bibliography{citation}

\end{document}





























